
\documentclass[11pt]{article}

\usepackage[a4paper,margin=1in]{geometry}
\usepackage{amsmath,amssymb,amsthm,mathtools,mathrsfs,bm}
\usepackage{enumitem}
\usepackage{hyperref}
\usepackage{amssymb}

\hypersetup{
    colorlinks=true,
    linkcolor=blue,
    citecolor=blue,
    urlcolor=blue
}

\title{
Exact Renewal Structure in Accelerated Collatz Dynamics\\
\large A 2-adic Symbolic Model and Exact Operator Theory
}

\author{Hiroki Kamanoi}
\date{2026}

\newtheorem{theorem}{Theorem}[section]
\newtheorem{lemma}[theorem]{Lemma}
\newtheorem{corollary}[theorem]{Corollary}
\newtheorem{remark}[theorem]{Remark}
\newtheorem{definition}[theorem]{Definition}

\begin{document}

\maketitle

\begin{abstract}
We construct an exact symbolic dynamical model associated with the accelerated
Collatz first-return map on the Type A congruence class
\[
\mathcal A := \{n\in\mathbb Z_{>0}: n\equiv3\pmod4\}.
\]

The paper separates the problem into two logically distinct layers.

First, we define an autonomous countable-state renewal system
\[
(X,\sigma,\pi)
\]
and establish its exact operator structure. The associated Markov operator
admits an exact shift-and-kill decomposition on mean-zero observables:
\[
P^Nf=S^Nf.
\]

From this identity we derive exact formulas for covariance decay,
spectral contraction, entropy, and total mutual information memory.

Second, we show that accelerated Collatz dynamics embeds densely into this
symbolic system through a 2-adic coding map.

The present work does not prove the Collatz conjecture. Its purpose is to
isolate and rigorously characterize the symbolic and operator-theoretic
structure underlying the induced Type A dynamics.
\end{abstract}

\tableofcontents

\section{Introduction}

The accelerated Collatz map is defined by
\[
T(n)=
\begin{cases}
n/2 & (n \text{ even}),\\
(3n+1)/2^{v_2(3n+1)} & (n \text{ odd}).
\end{cases}
\]

Rather than studying arbitrary orbits directly, we focus on the induced
first-return dynamics on the congruence class
\[
\mathcal A=\{n\equiv3\pmod4\}.
\]

The key observation is that the valuation coordinate
\[
k(n):=v_2(n+1)-1
\]
evolves according to a renewal mechanism:
\begin{itemize}
\item deterministic descent for \(k\ge2\),
\item geometric renewal at \(k=1\).
\end{itemize}

The main contribution of this paper is the separation between:
\begin{enumerate}[label=(\roman*)]
\item an autonomous symbolic renewal system,
\item an arithmetic realization arising from Collatz dynamics.
\end{enumerate}

The symbolic system itself admits a completely explicit operator theory.

\subsection{Notation and Dynamical Conventions}

We distinguish carefully between:
\begin{enumerate}[label=(\arabic*)]
\item arithmetic Collatz dynamics,
\item the autonomous renewal system,
\item the 2-adic symbolic realization.
\end{enumerate}

This separation is maintained throughout the paper.

\begin{center}
\begin{tabular}{c|c}
Symbol & Meaning \\
\hline
\(T\) & accelerated Collatz map \\
\(T_{\mathcal A}\) & first-return map on Type A integers \\
\(X\) & renewal shift \\
\(\sigma\) & left shift \\
\(P\) & Markov operator \\
\(S\) & shift-and-kill operator \\
\(\pi\) & stationary distribution \\
\(k(n)\) & valuation coordinate \(v_2(n+1)-1\)
\end{tabular}
\end{center}

\subsection{Topological Structure}

The renewal shift \(X\) is endowed with the product topology induced from the
discrete topology on \(\mathbb N\).

The completion
\[
\widehat{\mathcal A}
=
\{x\in\mathbb Z_2:x\equiv3\pmod4\}
\]
inherits the standard 2-adic topology.

Finite symbolic words correspond to finite congruence conditions modulo powers
of \(2\).

\section{Arithmetic Setup}

Define
\[
\mathcal A=\{n\in\mathbb Z_{>0}:n\equiv3\pmod4\}.
\]

Let
\[
T_{\mathcal A}
\]
denote the first-return map to \(\mathcal A\).

For \(n\in\mathcal A\), define
\[
k(n)=v_2(n+1)-1.
\]

\section{The Autonomous Renewal System}

\subsection{State Space}

Define the renewal shift
\[
X\subset\mathbb N^{\mathbb N}
\]
generated by transitions
\[
j\to j-1
\qquad (j\ge2),
\]
and
\[
1\to m
\qquad (m\ge1).
\]

\subsection{Markov Kernel}

Define the transition kernel
\[
P(j,j-1)=1
\qquad (j\ge2),
\]
and
\[
P(1,m)=2^{-m}.
\]

\begin{theorem}[Stationary Measure]
The probability distribution
\[
\pi(m)=2^{-m}
\]
is stationary and unique.
\end{theorem}

\begin{proof}
For \(m\ge1\),
\[
(\pi P)(m)
=
\pi(m+1)+\pi(1)2^{-m}
=
2^{-m}.
\]

Irreducibility and positive recurrence imply uniqueness.
\end{proof}

\section{Operator Structure}

\subsection{Markov Operator}

Define
\[
(Pf)(j)=\sum_m P(j,m)f(m).
\]

Explicitly,
\[
(Pf)(j)=f(j-1)
\qquad (j\ge2),
\]
and
\[
(Pf)(1)=E_\pi[f].
\]

\subsection{Shift-and-Kill Operator}

Define
\[
(Sf)(j)=f(j-1)\mathbf1_{j\ge2}.
\]

\begin{theorem}[Shift-and-Kill Identity]
For every mean-zero function \(f\in L^2(\pi)\),
\[
P^Nf=S^Nf.
\]
\end{theorem}

\begin{proof}
For \(N=1\),
\[
(Pf)(j)=f(j-1)=(Sf)(j)
\qquad (j\ge2),
\]
while
\[
(Pf)(1)=E_\pi[f]=0=(Sf)(1).
\]

Since \(P\) preserves the mean-zero subspace, induction yields
\[
P^{N+1}f=S^{N+1}f.
\]
\end{proof}

\section{Exact Mixing and Spectral Theory}

\begin{theorem}[Exact Covariance]
Let
\[
g(k)=k-2.
\]
Then
\[
\operatorname{Cov}_\pi(k_0,k_N)=2^{1-N}.
\]
\end{theorem}

\begin{proof}
By the shift-and-kill identity,
\[
\operatorname{Cov}_\pi(k_0,k_N)
=
\langle g,S^Ng\rangle_\pi.
\]

Direct summation gives
\[
\sum_{j>N}(j-2)(j-N-2)2^{-j}=2^{1-N}.
\]
\end{proof}

\begin{theorem}[\(L^\infty\)-Mixing]
For bounded observables \(f,g\),
\[
|\operatorname{Cov}_\pi(f(k_0),g(k_N))|
\le
\|f\|_\infty\|g\|_\infty 2^{-N}.
\]
\end{theorem}

\begin{theorem}[\(L^2\) Spectral Decay]
For mean-zero \(f\in L^2(\pi)\),
\[
\|P^Nf\|_{L^2(\pi)}
=
2^{-N/2}\|f\|_{L^2(\pi)}.
\]
\end{theorem}

\section{Entropy and Information}

\begin{theorem}[Entropy]
The measure-theoretic entropy satisfies
\[
h(\pi,\sigma)=\log2.
\]
\end{theorem}

\begin{theorem}[Total Mutual Information]
\[
\sum_{N=1}^\infty I_\pi(k_0;k_N)=2\log2.
\]
\end{theorem}

\section{Arithmetic Realization}

Define
\[
\widehat{\mathcal A}
=
\{x\in\mathbb Z_2:x\equiv3\pmod4\}.
\]

\begin{theorem}[2-adic Symbolic Realization]
The coding map
\[
\widehat\Phi:\widehat{\mathcal A}\to X
\]
is a homeomorphism.

Positive integers inject densely into \(X\) through this realization.
\end{theorem}

\begin{proof}
We divide the proof into four parts.

\textbf{Injectivity.}
Finite symbolic prefixes determine congruence classes modulo powers of \(2\).
Hence two points with identical symbolic itineraries agree modulo \(2^N\) for every \(N\),
and therefore coincide in \(\mathbb Z_2\).

\textbf{Surjectivity.}
Given an admissible symbolic sequence
\[
(k_1,k_2,\ldots)\in X,
\]
the finite prefix \((k_1,\ldots,k_r)\) determines a nonempty congruence class modulo
\(2^{M_r}\), where
\[
M_r=k_1+\cdots+k_r+r.
\]

These congruence classes define nested 2-adic balls
\[
B_1\supset B_2\supset\cdots
\]
with diameters tending to zero.

Compactness and completeness of \(\mathbb Z_2\) imply
\[
\bigcap_r B_r
\]
contains a unique point.

\textbf{Continuity.}
Cylinder sets in \(X\) correspond exactly to finite congruence conditions modulo powers
of \(2\). Hence inverse images of cylinders are open in the 2-adic topology.

\textbf{Continuity of the inverse.}
Finite 2-adic congruence information determines finite symbolic prefixes, so the inverse
map is continuous.
\end{proof}

\section{Scope and Limitations}

This paper establishes:
\begin{itemize}
\item an exact autonomous renewal system,
\item its operator structure,
\item exact covariance and entropy formulas,
\item a 2-adic arithmetic realization.
\end{itemize}

The paper does not prove the Collatz conjecture.

In particular, no statement is made concerning:
\begin{itemize}
\item global convergence,
\item divergent orbit exclusion,
\item orbitwise ergodicity,
\item arithmetic decorrelation.
\end{itemize}

\section{Conclusion}

The induced Type A dynamics admits an exact symbolic renewal structure whose
operator theory is completely computable.

The remaining obstruction appears not operator-theoretic, but arithmetic:
whether an individual orbit can sustain persistent deviations from the renewal
background.

\appendix

\section{Numerical Sanity Checks}

Finite computations confirm:
\begin{itemize}
\item exact covariance values for small \(N\),
\item symbolic congruence reconstruction,
\item renewal transition statistics.
\end{itemize}

These computations are included only as consistency checks and are not used in any proof.

\section{Entropy and KL Computations}

The mutual information formula follows from the deterministic survival of
non-renewed trajectories:
\[
I_\pi(k_0;k_N)
=
2\log2\cdot2^{-N}.
\]

Summing the geometric series yields
\[
\sum_{N\ge1}I_\pi(k_0;k_N)=2\log2.
\]

\begin{thebibliography}{99}

\bibitem{tao}
T. Tao,
\emph{Almost all orbits of the Collatz map attain almost bounded values},
Forum of Mathematics, Pi (2019).

\bibitem{walters}
P. Walters,
\emph{An Introduction to Ergodic Theory},
Springer, 1982.

\bibitem{kitchens}
B. Kitchens,
\emph{Symbolic Dynamics},
Springer, 1998.

\bibitem{lindmarcus}
D. Lind and B. Marcus,
\emph{An Introduction to Symbolic Dynamics and Coding},
Cambridge University Press.

\bibitem{sarig}
O. Sarig,
\emph{Thermodynamic formalism for countable Markov shifts},
Ergodic Theory and Dynamical Systems.

\bibitem{aaronson}
J. Aaronson,
\emph{An Introduction to Infinite Ergodic Theory},
AMS.

\end{thebibliography}

\end{document}
